% introduction.tex
Chemische stoffen worden in regelgeving, vergunningverlening en databanken
op uiteenlopende manieren geïdentificeerd: via naam, CAS-nummer, EINECS-nummer
of een handelsbenaming. Dit leidt tot de centrale vraag van dit onderzoek:
\emph{is een stof altijd dezelfde stof?}

In de praktijk blijken stoffen onder verschillende benamingen of
registratienummers te worden aangevraagd, soms met afwijkende regelgevende
implicaties. Dit onderzoek brengt die discrepanties in kaart op basis van een
dataset van requesten op stoffenlijsten.

\subsection{Achtergrond}
% Beschrijf relevante regelgeving (REACH, CLP, etc.)

\subsection{Onderzoeksvraag}
% Formuleer de concrete onderzoeksvragen
